% =====================================================================
%  Chapter 10 — EARLY WORK III: the S31 audit
% =====================================================================
\chapter{Early Work III: The Audit}
\label{ch:audit}

\begin{record}
\textbf{Classification.} Early work. Session \Sn{31}, 18 July 2026. The audit
produced no new mathematics and is the most consequential session in the
programme's history. Its findings are recorded in the proof document at
\texttt{sec:audit-s31}, line~5218.
\end{record}

The session began as something else. The owner opened it asking to explore a
new possibility for proving global existence. Before selecting a new route,
the programme performed an independent line-by-line audit of what it already
believed --- the vorticity scale recursion of \S20 and the geometric
framework of \S21--\S22 --- because those were the basis of a claim about to
be posted publicly. The audit found the claim invalid.

\section{What we assumed}
\label{sec:audit-assumed}

An audit's hypotheses are about the reading, not about the fluid. Four were in
force.

\begin{enumerate}[leftmargin=2.2em]
\item \emph{The document is read as written, not as intended.} Where a step is
  asserted, it is checked as an assertion. Where a constant is chosen, the
  choice is substituted back and the resulting inequality evaluated. No credit
  is given for a step that would work if a different constant had been chosen.
\item \emph{The quantifiers are taken seriously.} \S20's theorem quantifies
  over Leray--Hopf solutions with \(u_0\in H^1(\T^3)\). Every step must
  therefore be legitimate for a Leray--Hopf solution, not merely for a smooth
  one. This single discipline produced Defect~D4.
\item \emph{Every exponent is checked arithmetically.} H\"older triples are
  summed; Jensen's inequality is checked for direction; comparison ODEs are
  integrated.
\item \emph{The audit does not repair.} Its output is a list of defects and a
  revised status, not a corrected proof. Repair, where possible, is separate
  work under separate scrutiny.
\end{enumerate}

The last of these is a real methodological choice with a cost. It means the
audit leaves the document in a worse state than it found it --- claims
withdrawn, nothing put in their place --- and that state is what the programme
then had to work from. The alternative, an audit that simultaneously repairs,
would have blurred the boundary between what was found broken and what was
subsequently believed fixed.

\section{What we considered}
\label{sec:audit-considered}

\subsection{Repair or reformulate}

Once the defects were established, the question was what to do with \S20. Two
options were weighed.

\emph{Repair.} Re-derive the stretching exponent with valid H\"older and
Jensen steps, inside a framework where the local enstrophy inequality is
legitimate. This is bounded, well-defined work. The audit's own assessment of
the likely outcome is recorded: ``expected honest outcome is a conditional
(small-data / CKN-type) statement'' --- that is, a repaired \S20 would
reproduce what Caffarelli--Kohn--Nirenberg already gives, and no more. Nobody
has attempted it.

\emph{Reformulate.} Take what survived and build something new on it. This is
what happened, and \S\ref{sec:audit-produced} explains why the survivors made
it attractive.

\subsection{Four candidate routes}

The audit closed with a feasibility brainstorm over four directions. They are
reproduced here with the reasoning that was given at the time, because three
of the four were rejected on grounds that still hold.

\begin{center}
\small
\begin{tabular}{@{}l p{0.50\textwidth} l@{}}
\toprule
Route & Assessment at the time & Chosen? \\
\midrule
Alignment pincer
  & Most promising. The gap to bridge --- Theorem-E-style alignment is
    \(L^2\)-averaged, while Constantin--Fefferman needs pointwise or local
    direction coherence --- is a genuine open research question, \emph{not} a
    restatement of Claim~A. Prerequisite: repair Defects~D5--D8.
  & \textbf{yes} \\[3pt]
De~Giorgi--Nash--Moser on \(\theta\)
  & Obstruction identified up front: the source \(\nu\abs{\nabla u}^2\) is in
    \(L^1\) only, scalar theory then gives \(\theta\in L^q\) for \(q<5/3\) and
    no H\"older gain. Any improvement must exploit that the source is the
    dissipation of the drift's own energy --- which is Claim~A restated.
    Worth one bounded session, documented as such.
  & no \\[3pt]
Probabilistic
  & Almost-sure global well-posedness for randomised \(H^1\) data. Publishable,
    but a different theorem from the Prize.
  & no \\[3pt]
\S20 repair
  & See above; expected to land in the CKN small-data regime.
  & no \\
\bottomrule
\end{tabular}
\end{center}

\medskip
The reason the first was chosen is worth stating precisely, because it is the
only one of the four that is not circular. The gap it must close --- upgrading
an averaged coherence bound to a pointwise one --- is a question about
parabolic regularity for the direction field, and the direction field's
equation is a genuinely different object from the Navier--Stokes equations
themselves. It is not the Prize problem wearing a hat. Everything in Part~IV
of this book follows from that judgement.

\section{What it produced}
\label{sec:audit-produced}

\subsection{Eight defects}

\begin{table}[ht]
\centering
\footnotesize
\begin{tabular}{@{}l p{0.55\textwidth} l@{}}
\toprule
\# & Defect & Kills \\
\midrule
D1 & \textbf{The induction is arithmetically impossible for large data.} The
  inductive step needs \(C_{\mathrm v}D^{5/3}2^{10/9}r^{10/9}\le\tfrac D2
  r^{2/3}\), i.e.\ \(D^{2/3}\le2^{-19/9}/C_{\mathrm v}\): \(D\) must be
  \emph{small}. The base case needs \(D\ge\norm{u_0}_{H^1}^2T\): \(D\) must be
  \emph{large}. The stated choice satisfies neither, and substituting it makes
  the superlinear term strictly larger than \(D/2\) whenever
  \(C_{\mathrm v}\ge1\). The recursion collapses to the CKN small-energy
  regime.
  & \texttt{thm:uniform-} \texttt{enstrophy} \\[3pt]
D2 & \textbf{Invalid H\"older triple.} \(\tfrac13+\tfrac12+\tfrac12=\tfrac43>1\).
  The correct pairing \(\norm{\nabla u}_{L^3}\norm{\omega}_{L^3}^2\) changes
  every downstream exponent; Young then produces
  \(\norm{\omega}_{L^2}^6\), not \(\norm{\omega}_{L^2}^{10/3}\). The claimed
  \(5/3\) superlinearity does not survive.
  & the recursion \\[3pt]
D3 & \textbf{Jensen reversed.} The step asserts
  \(\int g^{5/3}\le r^{-2/3}(\int g)^{5/3}\); Jensen gives the reverse.
  Repairing it needs \(L^\infty_t\) information, unavailable for Leray--Hopf
  solutions.
  & the recursion \\[3pt]
D4 & \textbf{Regularity circularity and norm mixing.} Testing the vorticity
  equation against \(\varphi_r^2\omega\) is justified for smooth solutions but
  not for Leray--Hopf ones: \(\nabla\omega\in L^2_{\mathrm{loc}}\) is not known
  a priori, and suitable weak solutions satisfy a local \emph{energy}
  inequality, not a local \emph{enstrophy} inequality. Separately, global
  norms are used where local ones are required.
  & the local lemma \\[3pt]
D5 & \textbf{Conditional on Beale--Kato--Majda.} The hypothesis
  \(\int_0^TM\,dt<\infty\) \emph{is} the BKM criterion; under it regularity
  already follows classically. The theorem assumes an equivalent of its
  conclusion.
  & \texttt{thm:sigma-} \texttt{gronwall} \\[3pt]
D6 & \textbf{Large-viscosity condition.} The damping constant
  \(\beta=c\nu/C_1-C>0\) requires \(\nu>C_1C/c\), a relation between the given
  physical viscosity and universal constants. Navier--Stokes scaling cannot
  arrange it.
  & Regime~I \\[3pt]
D7 & \textbf{The Regime~II conclusion is false.} From
  \(\dot M\le CM^{5/4}(\log M)^{1/2}\) the text concludes \(M\) ``stays finite
  on any fixed time interval''. The comparison ODE \(\dot M=CM^{5/4}\) has
  solution \(M(t)=(M(0)^{-1/4}-\tfrac C4t)^{-4}\), which blows up at
  \(t^*=4M(0)^{-1/4}/C\); the logarithm only accelerates it.
  & Regime~II \\[3pt]
D8 & \textbf{Unproved coercivity, and the \(L^\infty\)
  Calder\'on--Zygmund bound.} The claim that diffusion contributes
  \(-c\nu\int\abs{\nabla^2\omega}^2\varphi^2\) to
  \(\tfrac{d}{dt}A_{\mathrm{loc}}\) is asserted without proof: differentiating
  \(A_{\mathrm{loc}}\) involves \(\ehat=\omega/\abs{\omega}\), singular on
  \(\{\omega=0\}\), and produces sign-indefinite cross terms. Moreover the
  stretching estimate invokes
  \(\norm{\nabla u}_{L^\infty(B_{\rstar})}\le CM\), \emph{but
  Calder\'on--Zygmund operators are unbounded on \(L^\infty\), and the correct
  bound carries a logarithmic factor} --- which matters precisely because
  \S22 relies on log-corrected estimates.
  & \texttt{lem:A-evol}, and much more \\
\bottomrule
\end{tabular}
\caption[The eight defects]{The eight grounds of the \Sn{31} audit. Note the
count: \emph{eight}, not seven. \texttt{PROGRESS.md} lists seven, omitting D8;
the proof document governs, and \Sn{47} corrected the summaries
accordingly.}
\label{tab:eight-defects}
\end{table}

\begin{figure}[ht]
\centering
\begin{tikzpicture}[
  >={Latex[length=1.8mm]},
  n/.style={draw, rounded corners=2pt, font=\scriptsize, align=center,
            inner sep=3pt, minimum height=8mm, text width=24mm},
  ok/.style={n, draw=provedgreen, thick},
  bad/.style={n, draw=impossiblered, very thick},
  d/.style={circle, draw=impossiblered, fill=impossiblered!12,
            font=\scriptsize\bfseries, inner sep=1.2pt, minimum size=6.5mm},
  ar/.style={->, thick}
]
% S20 chain
\node[font=\small\bfseries, anchor=west] at (-2.4,1.0) {\S20};
\node[ok]  (h1) at (0,1.0)  {\(u_0\in H^1\)};
\node[bad] (rec) at (3.2,1.0) {recursion\\\(Z(r)\le CZ(2r)^{5/3}\)};
\node[bad] (unif) at (6.6,1.0) {\(Z(r)\le D\,r^{2/3}\)};
\node[bad] (prize) at (10.0,1.0) {\textbf{Prize for \(H^1\)}};
\draw[ar] (h1)--(rec); \draw[ar] (rec)--(unif); \draw[ar] (unif)--(prize);

\node[d] (dd2) at (2.5,2.5) {D2};
\node[d] (dd3) at (3.9,2.5) {D3};
\node[d] (dd4) at (3.2,-0.3) {D4};
\node[d] (dd1) at (6.6,2.5) {D1};
\draw[ar, impossiblered] (dd2)--(rec.north west);
\draw[ar, impossiblered] (dd3)--(rec.north east);
\draw[ar, impossiblered] (dd4)--(rec.south);
\draw[ar, impossiblered] (dd1)--(unif);

% S21-22 chain
\node[font=\small\bfseries, anchor=west] at (-2.4,-2.2) {\S21--\S22};
\node[ok]  (split) at (0,-2.2) {split identity\\\(\sigma\) invariance};
\node[bad] (aevol) at (3.2,-2.2) {\texttt{lem:A-evol}\\coercivity};
\node[bad] (gron) at (6.6,-2.2) {\texttt{thm:sigma-}\\\texttt{gronwall}};
\node[bad] (thmE) at (10.0,-2.2) {Theorem E\\alignment};
\draw[ar] (split)--(aevol); \draw[ar] (aevol)--(gron); \draw[ar] (gron)--(thmE);

\node[d] (dd8) at (3.2,-3.9) {D8};
\node[d] (dd5) at (6.0,-3.9) {D5};
\node[d] (dd6) at (7.4,-3.9) {D6};
\node[d] (dd7) at (10.0,-3.9) {D7};
\draw[ar, impossiblered] (dd8)--(aevol);
\draw[ar, impossiblered] (dd5)--(gron.south west);
\draw[ar, impossiblered] (dd6)--(gron.south east);
\draw[ar, impossiblered] (dd7)--(thmE);

\draw[->, thick, openblue, dashed] (dd8) to[out=180,in=200,looseness=1.3]
  node[left, font=\scriptsize, text=openblue, align=right, pos=0.55]
  {and forward,\\to every later\\session} (split.west);
\end{tikzpicture}
\caption[Where the eight defects strike]{The two chains and where each defect
strikes. Green boxes are the survivors. The dashed blue arrow is the subject
of Chapter~\ref{ch:inheritance}: D8's second half --- the \(L^\infty\)
Calder\'on--Zygmund bound --- is not merely a defect in \texttt{lem:A-evol}
but the first appearance of the obstruction that dominates \Sn{32}--\Sn{47}.}
\label{fig:audit-defects}
\end{figure}

\subsection{The survivors}

The audit's own list, at \texttt{sec:audit-s31}, of what remains intact:

\begin{enumerate}[leftmargin=2.2em]
\item The Route~C bootstrap, \texttt{thm:routeC} --- unconditional. Intact.
\item Claim~A for shear-layer initial data (\S11) and near-shear data (\S12).
  Intact.
\item \(\lambda_{\max}=0\) for shear flows. Intact.
\item The algebraic identity \texttt{lem:grad-omega-split} and the exact
  scale invariance of \(\sigma=A_{\mathrm{loc}}/M^{3/2}\). Intact ---
  ``sound foundations for a future quantitative alignment program''.
\item \S20 as a \emph{conditional} small-data statement, consistent with CKN,
  \emph{provided} D2--D4 are repaired and the superlinearity exponent
  re-derived.
\end{enumerate}

Everything else is withdrawn or conditional. Note what is \emph{not} on the
list: Theorem~E is absent, as are \texttt{prop:coercive},
\texttt{thm:sigma-gronwall} and the log-corrected BKM criterion.

\subsection{Revised status, and the actions taken}

\begin{record}
\textbf{Revised status (\texttt{sec:audit-s31}).} Global regularity for
general \(u_0\in H^1(\T^3)\) is \OPEN{} in this document. The Prize claim of
\S20 (\texttt{cor:prize-h1}) is \WITHDRAWN. The programme's strongest verified
position: unconditional regularity for shear and near-shear initial data via
Claim~A and Route~C, plus extensive numerical evidence consistent with --- but
not proving --- the general case.
\end{record}

The propagation was done the same day and is worth listing, because a
withdrawal that is not propagated is not a withdrawal: the audit section was
appended to the proof document and its status table updated; \texttt{PROGRESS.md}
was downgraded; the preprint's abstract was downgraded and its main theorem
demoted to a conjecture with a withdrawal remark; the programme's overview
document was downgraded; and \texttt{arXiv} submission was halted.

A test-count discrepancy surfaced in the same pass and was never resolved:
the audit measured \(904\) passing tests against the \(873\) recorded in
\texttt{PROGRESS.md} at \Sn{30}, and the proof document's own revised-status
paragraph quotes \(873\). The delta appears to be test additions never counted
into the summary. No failures anywhere.\footnote{This is a live source
disagreement of exactly the kind \S\ref{sec:source-hierarchy} anticipates.
The proof document says \(873\); the \Sn{31} session log says \(904\)
observed; by \Sn{45} the \texttt{layer4} suite alone reports \(443\). This
book quotes the document's figure where the document is the source, and the
logs' figures where a session log is the source, and does not attempt to
reconcile them.}

\section{What survived and what did not}
\label{sec:audit-survived}

The audit is not a proof, so the usual status labels do not quite apply to it.
What can be asked is whether its findings have held.

\begin{center}
\small
\begin{tabular}{@{}p{0.32\textwidth} p{0.15\textwidth} p{0.45\textwidth}@{}}
\toprule
Finding & Status today & Note \\
\midrule
The withdrawal of \texttt{cor:prize-h1} and the two corollaries
  & \WITHDRAWN, permanently
  & Reaffirmed as Wall~0 in \Sn{47}; nobody has attempted the repair \\[3pt]
D1--D4 as grounds
  & Standing & Never contested \\[3pt]
D5--D7 as grounds
  & Standing & Never contested \\[3pt]
D8, first half (coercivity unproved)
  & Standing & \texttt{lem:A-evol} was never repaired \\[3pt]
D8, second half (the \(L^\infty\) CZ bound carries a logarithm)
  & \textbf{Standing, and structural}
  & Became \(\Lfac\); rediscovered three further times; consolidated as
    Wall~1. See Chapter~\ref{ch:inheritance} \\[3pt]
The five-item survivors list
  & All five still hold & Items 1--4 are still load-bearing; item~5 is
    dormant \\[3pt]
The four-route feasibility assessment
  & Route~1 chosen; the three rejections still stand
  & The alignment pincer produced Part~IV \\
\bottomrule
\end{tabular}
\end{center}

\begin{obsv}[The audit's most valuable finding was the one nearly lost]
\label{obs:d8-lost}
Of the eight defects, seven are local: each kills a specific step, and each is
in principle repairable by redoing that step correctly. The eighth is not.
D8's second half identifies a \emph{structural} feature of the whole
architecture --- that recovering \(\nabla u\) from \(\omega\) by Biot--Savart
and then taking a supremum costs a logarithm, unavoidably, because
Calder\'on--Zygmund operators are unbounded on \(L^\infty\). Every estimate in
the programme that follows this architecture inherits it.

And D8 is the defect that the programme's own summary dropped. For sixteen
sessions \texttt{PROGRESS.md} recorded the audit as having seven grounds. The
consequence was not merely a miscount: when the logarithm reappeared in
\Sn{34}, and again in \Sn{41}, and again in \Sn{44}, it looked each time like
a fresh difficulty rather than the same one, and there was no record saying
otherwise. The \Sn{47} walls document restores the connection and states it
explicitly --- the factor ``is not a later regression but the audit's finding
carried forward''.
\end{obsv}

\begin{obsv}[What an audit costs and what it buys]
\label{obs:audit-economics}
The audit cost the programme its headline result and ten weeks of preparation
for a submission that never happened. What it bought is the only reason this
book can be read at all: every claim after \Sn{31} was made in a programme
that had already destroyed one of its own, and knew it could happen again.
The verification discipline of Chapter~\ref{ch:house-rules} is entirely
downstream of this session, and so, in a more literal sense, is Part~IV --- the
alignment pincer exists because the audit both cleared the ground and left
exactly two sound objects standing on it.
\end{obsv}
